% RESUMEN

\VIUSectionStar{Resumen}

Este Trabajo Fin de Máster desarrolla un módulo para monitoreo de calidad del aire orientado a la evaluación del riesgo y a la generación de alertas en entornos urbanos. La propuesta parte de una revisión sistemática de la literatura sobre sistemas de inferencia difusa aplicados a monitoreo en tiempo real, con el fin de identificar variables dominantes, enfoques difusos recurrentes, arquitecturas operativas y vacíos metodológicos que condicionan el diseño de la solución. A partir de esa evidencia se definió un artefacto híbrido que combina una capa normativa de cálculo del índice de calidad del aire, una capa de inferencia difusa principal orientada a interpretar concurrencia y persistencia, y una capa contextual desacoplada basada en reglas para temperatura y humedad.

La implementación se materializó en un prototipo ejecutable denominado \textit{AQRisk}, construido en Python mediante una arquitectura modular. El sistema integra ingesta de datos abiertos, preparación y control de calidad, cálculo del AQI con base en \textit{AQI Breakpoints} de \textit{EPA/AQS}, un motor Mamdani con una base principal de 54 reglas, una base contextual crisp de 9 reglas y un módulo de alertamiento trazable. Como complemento, se desarrolló una interfaz web para control de entrada, visualización de resultados, explicabilidad del razonamiento y revisión histórica de ejecuciones.

Los resultados actuales muestran que el prototipo ejecuta el flujo completo del artefacto, conserva trazabilidad por capas, reglas y variables derivadas, y permite validar escenarios controlados y consultas reales sobre datos abiertos. El aporte principal del trabajo consiste en articular, dentro de un mismo marco reproducible, cálculo normativo, inferencia difusa explicable, alertamiento interpretable e interfaz de supervisión para apoyo a la decisión.

\textbf{Palabras clave:} calidad del aire, lógica difusa, AQI, monitoreo en tiempo real, alertamiento, datos abiertos, explicabilidad.

\clearpage
\VIUSectionStar{Abstract}

This Master's Thesis develops an air-quality monitoring module oriented to risk assessment and alert generation in urban environments. The proposal is grounded on a systematic literature review of fuzzy inference systems applied to real-time monitoring in order to identify dominant variables, recurring fuzzy approaches, operational architectures, and methodological gaps that shape the design of the solution. Based on that evidence, a hybrid artefact was defined, combining a normative air-quality-index layer, a main fuzzy inference layer aimed at interpreting pollutant concurrence and temporal persistence, and a decoupled rule-based contextual layer for temperature and humidity.

The implementation was materialized in an executable prototype named \textit{AQRisk}, built in Python through a modular architecture. The system integrates open-data ingestion, data preparation and quality control, AQI computation based on the \textit{EPA/AQS AQI Breakpoints}, a Mamdani engine with a 54-rule main base, a 9-rule crisp contextual base, and a traceable alerting module. In addition, a web interface was developed to control inputs, visualize outputs, explain the reasoning process, and review execution history.

Current results show that the prototype executes the complete pipeline, preserves traceability by layers, rules, and derived variables, and supports both controlled scenarios and real open-data queries. The main contribution of the work lies in articulating, within a single reproducible framework, normative computation, explainable fuzzy inference, interpretable alerting, and a supervision interface for decision support.

\textbf{Keywords:} air quality, fuzzy logic, AQI, real-time monitoring, alerting, open data, explainability.
